\label{texfile:GWF}
\includegraphics[width=.25\textwidth]{UnderConstruction} \textit{
Should add some into material about GWF build}.

\subsection{Generating a Layered \gwf\ Domain}
A \mfus\ 3D groundwater flow (\gwf) domain can be generated from the template using this instruction:

\ins{generate layered gwf domain}
    {This subtask has instructions that are used to define:
     \begin{itemize}
        \item Element zone numbering scheme
        \item Top elevation (i.e. $z$-coordinate)
        \item Mesh layers and vertical discretization
    \end{itemize}

    Subtask instructions will be read and processed until an \textsf{end} instruction is encountered.  We suggest appending the subtask name to the \textsf{end} instruction:

    {\Large \sf end generate layered gwf domain}
    }

The construction of the 3D \gwf\ finite-element mesh proceeds from top to bottom.  First, we define the top elevation, then add layers one at a time until we reach the base of the domain. By default, element zone numbers will be assigned by layer number.

\subsubsection{Defining the Top Elevation}
To assign an elevation to the top layer of template nodes use this instruction:

\ins{top elevation}
    {This subtask defines the elevation (i.e. $z$-coordinate) of the top layer of nodes in the \gwf\ finite-element template mesh in one of these ways:
     \begin{itemize}
        \item By assigning a given elevation to all nodes
        \item By reading variable elevation data from a file
        \item By interpolating elevation data from a function $z(x)$ where the elevation $z$ varies by the nodes $x$ coordinate.
     \end{itemize}
     One the elevation is defined, an \textsf{end} instruction is required to stop the subtask e.g.\:

    {\Large \sf end top elevation}
    }

 The top elevation can be defined by one of these instructions: \label{'Page:TopElev'}

 \ins{elevation constant}
    {\squish
    \begin{enumerate}
    \item \rnum{elev}\  The elevation \rnum{elev}\ will be assigned to all top layer nodes.
    \end{enumerate}
    \squish
    }

 \ins{elevation from gb file}
    {\squish
    \begin{enumerate}
    \item \str{file}\  The elevation data in the \gb\ nodal property file named \str{file}\ will be assigned to the top layer nodes.
    \end{enumerate}
    \squish
    }

The \gb\ nodal property file uses a legacy binary file format. You can develop your own ascii input files and read them using this instruction:

 \ins{elevation from list file}
    {\squish
    \begin{enumerate}
    \item \str{file}\  The elevation data in the ascii file named  \str{file}\ will be assigned to the top layer nodes.
    \end{enumerate}
    \squish
    }

Part of a sample list file \footnote{The verification example \texttt{MUT\_Examples$\backslash$6\_Abdul\_MODHMS} uses an ascii file input to define nodal elevations.} is shown here:
    \begin{verbatim}
    Kriged cell top elevation for layer 1
     4.414571762E+000
     4.415914536E+000
     ...
     4.415914536E+000
     \end{verbatim}
     \squish
Some key features of this example are:
\begin{itemize}
  \item The first line of the file is discarded, and in this case contains a string describing the data.
  \item You must supply a value for each node in the template finite-element mesh.
  \item The data is read in free format so there can be more than one value entered per line.
  \item Only the start and end of the file are shown here, with the string '\texttt{...}' replacing the middle portion.
\end{itemize}

To define the top elevation as a function of $x$ (usually used for cross-sectional models) use this instruction:

 \ins{elevation from xz pairs}
    {
    \squish
    \begin{enumerate}
    \item \rnum{x(1)}, \rnum{y(1)}  First $x, y$ coordinate pair.
    \item \textbf{...}
    \item \rnum{x(n)}, \rnum{y(n)}  nth $x, y$ coordinate pair.
    \end{enumerate}

     This subtask reads a list of $xz$-coordinate pairs until an \textsf{end} instruction is encountered e.g.\:

    {\Large \sf end elevation from xz pairs}
    }

A sample showing the use of this instruction \footnote{The verification example \texttt{MUT\_Examples$\backslash$1\_VSF\_Hillslope} uses the $xz$ pairs instruction to define the top elevation of the cross-sectional domnain.} is shown here:
    \begin{verbatim}
    elevation from xz pairs
           0.0, 0.0
        1000.0, 100.0
    end elevation from xz pairs
     \end{verbatim}
     \squish
Some key features of this example are:
\begin{itemize}
  \item The two given $xz$ pairs define a line that slopes from $z=0$ at $x=0$ to $z=100.0$ at $x=1000$.  You may supply as many pairs as needed to define the top of your cross-section.
  \item $x$ coordinates must increase continuously from the top of the list to the bottom.
  \item the $x$-range of the supplied pairs should cover the entire $x$-range of the template mesh.
  \item For each node in the template mesh, the $x$ coordinate is used to interpolate an elevation (i.e. $z$ value) using the appropriate $xz$ pair.
\end{itemize}

\subsubsection{Adding Layers} \index{\gwf\ Domain ! layering}
    {\em NOTE: The term layers used here should not be confused with the \mf\ term of the same name. A \mf\ layer is one cell thick, while a \mut\ layer can be one or more elements thick.}

A \mfus\ model must contain at least 1 layer, and each layer is defined using this instruction:

\ins{new layer}
    {\index{Grid generation ! \gwf\ Domain ! New layer}
    This subtask adds a new layer to the \gwf\ domain by defining the layer:
     \begin{itemize}
       \item Base elevation
       \item Vertical discretization
     \end{itemize}

     It reads instructions until an \textsf{end} instruction is found e.g.\:

    {\Large \sf end new layer}
    }

The base elevation is defined using the elevation instructions  described on page~\pageref{'Page:TopElev'} that are given for the \textsf{top elevation} instruction.

Be default, a new layer will be assigned the name '\texttt{Layer {\em n}}' where {\em n} is the current layer number.  If you want to assign your own layer name use this instruction:

\ins{Layer name}
    { \squish
    \begin{enumerate}
    \item \str{layer\_name} Layer name.
    \end{enumerate}
    Assigns the layer name \str{layer\_name}.
    }

\includegraphics[width=.15\textwidth]{UnderConstruction} \textit{These names are not currently used in Tecplot output but could/should? be used to create the customlables for zone naming.}
    
    

By default, \mut\ will stop and issue a warning message if the computed layer base elevation is greater than or equal to the current layer top elevation.  This instruction forces the base to be below the top by a set amount:

\ins{Minimum layer thickness}
    {\squish
    \begin{enumerate}
    \item \rnum{MinThick} Minumum thickness value[L].
    \end{enumerate}
    This instruction causes \mut\ to enforce a minimum thickness constraint for the current layer. At nodes where the computed layer base elevation is greater than or equal to the current top elevation, \rnum{MinThick} will
    be subtracted from the current top elevation to get the base elevation.
    }

By default, a new layer will not be subdivided vertically unless one the following two instructions is issued.  
The first creates a uniform subdivision:

\ins{Uniform sublayering}
    {\squish
    \begin{enumerate}
    \item \inum{nsublayer} Number of sublayers.
    \end{enumerate}
    This instruction divides the layer vertically into \inum{nsublayer}
    elements, which will each have the same element height, equal to the top elevation
    minus the current base elevation divided by \inum{nsublayer}.
    }

This instruction creates a non-uniform subdivision:

\ins{Proportional sublayering}
    {\squish
    \begin{enumerate}
        \item \inum{nsublayer}  Number of proportional sublayers.
        \item \rnum{sub\_thick(i),i=1,\inum{nsublayer}} Proportional thicknesses in order from top to bottom.
    \end{enumerate}
    This instruction can be used if you want to refine the \gwf\ domain mesh vertically,
    for example, in the active zone with the \swf\ domain the ground surface in the .

    It is important to understand that the variable \rnum{sub\_thick} is not
    a true thickness, but is instead a relative thickness, which is used along
    with the layer thickness to determine the element heights in the current
    column.

}

    For example, these instructions:  
    \texttt{
            \begin{tabbing}
            AAA\=AAA\=AAA   \kill
                \>Proportional sublayering   \\
                \>   \> 3                     \\
                \>   \> 0.1           \\
                \>   \>1.0            \\
                \>   \>  10.0       \\
                \> end       \\
            \end{tabbing}
    }
    \squish
would subdivide the current layer vertically into three elements, between
    the current base and top elevation, with element height proportions of .1,
    1 and 10 from top to bottom.

This instruction is most often used to define a layer of uniform thickness relative to an uneven top elevation: 

\ins{Offset base}
    {\squish
    \begin{enumerate}
    \item \rnum{value} Thickness value (L) by which to offset the layer base elevation.
    \end{enumerate}
    This instruction causes the elevation of the base of the layer to be offset vertically by the given value.  This
    can be used to create a surface a given distance below another surface.
}

    For example, these instructions:
\begin{verbatim}
        top elevation
            elevation from list file
            elev.list
        end top elevation

        new layer
            uniform sublayering
            3

            elevation from list file
            elev.list

            offset base                     
            -1.0                           

        end new layer

    end generate layered gwf domain
\end{verbatim}
 create a layer with a top elevation 1 metre below the elevation defined in the raster file \texttt{gs.asc}:
 
 By default, element zone numbers reflect the horizontal layers added during mesh generation.
 If the template mesh is divided into horizontal patches with unique zone numbers, these can be assigned instead to the 3D GWF mesh \footnote{The verification example \texttt{MUT\_Examples$\backslash$6\_Abdul\_Prism\_Cell} uses this option to define \swf\ domain zones.} using this instruction:

\ins{Zone by template}
    {Causes \mut\ to assign the template mesh element zone number to the corresponding 3D \gwf\ element.

    {\em This instruction should appear in the input file at the beginning of the \textsf{generate layered gwf domain}}
    }

\begin{verbatim}
    generate layered gwf domain

        top elevation
            elevation constant
            100.0
        end top elevation

        new layer
            layer name
            Top layer

            uniform sublayering
            100

            elevation constant
            0.0
        end new layer

    end generate layered gwf domain
\end{verbatim}
The \textsf{generate layered gwf domain} instruction is a subtask so it has a corresponding \textsf{end generate layered gwf domain} instruction.  The \gwf\ domain is generated from the top layer down, and the first instruction, \textsf{top elevation}, defines the ground surface elevation at the layer of elements in the TECPLOT\_GWF mesh. It must begin with a \textsf{top elevation} instruction and at least one \textsf{new layer} instruction. Not shown in this example is an optional additional instruction, \textsf{zone by template}.
